% Expanded replacement for Discussion and Conclusion/Future Work.
% Replace the former Parts VI and VII in their entirety with these sections.

\section{Discussion}

The principal outcome of this work is the replacement of implicit handoffs with explicit scenario artifacts. In the earlier ECLIPSE demonstration, a user could define a lunar-base architecture, execute a DES, and inspect a visualization, but the operational assumptions and visualization state connecting these stages were partly hidden in separate scripts or prepared assets. The revised workflow makes the scenario itself the object that moves through the ecosystem. Its architecture, placement, routes, equations, numerical assumptions, results, CAD associations, and Omniverse files are preserved under a named scenario. This does not remove every implementation dependency, but it provides a traceable path from a user decision to the quantitative and visual products used to assess that decision.

\subsection{A. Automation of the Simulation-to-Visualization Link}

The scenario-specific manifest is the central automation mechanism. Rather than asking the Omniverse extension to reconstruct the mission from a fixed internal scenario, the upstream workflow writes the scene, waypoint, terrain, actor, and DES information required for one completed scenario. The extension selects and reads that package. This separation is important because it gives the visualization a clear scope: it should render and inspect the scenario produced by the engineering workflow, rather than become another location where scenario logic is manually reimplemented.

The manifest also improves reproducibility. Independent scenario packages prevent the most recent DES run from silently replacing the files used by a previous visualization. This allows a user to return to a saved alternative, inspect the same routes and CAD assets, and associate the playback with the corresponding result set. The approach remains dependent on the present GitHub-based execution architecture and a local repository pull within Omniverse. It should therefore be understood as automated pipeline integration, not as a continuously deployed cloud visualization service.

Terrain grounding and route projection provide an additional level of coherence. They reduce the discrepancy between the two-dimensional mission definition and the three-dimensional terrain used for inspection. However, this coherence is conditional on the supplied terrain mesh. The current workflow aligns and samples the prepared terrain asset associated with a scenario; it does not yet automatically retrieve, crop, and preprocess the exact terrain patch for every possible site selected in the interface. The slope overlay is consequently valuable not only as a visual feature but also as a diagnostic for potential inconsistency between the terrain representation and upstream route-feasibility analysis.

\subsection{B. From a Configurable ISRU Model to an Extensible Behavior Framework}

The DES work has two distinct implications. First, the ISRU model is no longer an opaque controller with a few exposed sliders. The user can inspect and modify the architecture, routes, equations, fleet configuration, production settings, storage settings, and power assumptions that govern the scenario. This is a substantial improvement in interpretability: when a result changes, the user can identify the scenario decision responsible for the change. The saved configuration also permits the same operational assumptions to be rerun or compared later.

Second, the class-based framework provides a way to generalize this process. The important contribution is not a claim that a generic engine automatically understands every subsystem submitted in SysML. Instead, it defines a limited operational language in which a user and the pipeline can describe a mission through sources, processors, storage elements, consumers, transporters, power generators, routes, and constrained equations. This language turns graph structure and interface information into executable SimPy behavior. The cargo-relay, ice-processing, and water-transfer examples demonstrate that the workflow can represent resource chains other than the original regolith-to-LOX case.

This generalization also establishes a useful boundary between model configuration and model development. A user can assemble scenarios that fit the available roles and equation contracts without creating a new Python controller. A new physical phenomenon that does not fit those roles, such as complex thermal behavior, stochastic equipment degradation, or coupled multi-resource chemistry, still requires an extension of the behavior library and validation of the corresponding model. The framework therefore improves flexibility while making the remaining modeling assumptions explicit.

\subsection{C. CAD Metadata and Authoritative-Source Consistency}

The CAD workflow addresses another common source of incoherence in digital engineering environments: geometry is often treated as a visualization attachment that is disconnected from the engineering model. In ECLIPSE, the CAD submission now produces both a standardized USD asset and traceable metadata products. The separation between source-CAD metadata and visualization metadata preserves the distinction between what was delivered by a stakeholder and what was generated to render it in Omniverse. This is particularly important for converted mesh formats, where source properties, units, orientation conventions, and physical calculations may not be represented identically after conversion.

Maintaining SysML as the dimensional authority avoids a competing source of truth. The imported CAD asset can be visually adapted to the architectural dimensions without redefining them, while scaled area, volume, and center-of-mass information can still enrich the SysML model. The approach is suitable for the current prototype, but physical metadata should be interpreted according to its provenance and reliability. In particular, volume is only meaningful when the source mesh is sufficiently watertight, and a CAD model should not be assumed to represent an engineering-certified mass model merely because it produces a numerical center of mass or surface area.

\subsection{D. Current Limitations}

Several limitations define the current boundary of the contribution. The generic DES supports a constrained role vocabulary and equation language; it does not yet provide comprehensive unit checking, detailed loading and unloading equipment, rich queues, stochastic events, complex multi-resource synchronization, or a high-fidelity electrical network. The power model is sufficient to expose generation, storage, and consumption assumptions and to detect modeled energy shortfalls, but it is not a detailed electrical-system design tool. Similarly, rover movement in Omniverse follows terrain-projected routes, whereas the DES uses route-level travel metrics rather than full vehicle dynamics.

The time-synchronized Omniverse dashboard should also be interpreted carefully. It reads the completed DES log at the selected playback time, which is appropriate for reproducible inspection of a result. It does not yet update while a cloud-hosted DES is executing. The saved-scenario comparison capability is likewise an initial design-space exploration tool; its available histories and aggregated metrics should become more resource-agnostic as the generic engine evolves. Finally, the present use of GitHub Pages, GitHub Actions, and local Omniverse repositories is practical for a prototype, but it cannot provide the transactional persistence, authorization, conflict resolution, and simultaneous editing required for true real-time multi-user mission design.

\section{Conclusions and Future Work}

This work advances ECLIPSE from a manually coupled lunar-base demonstration toward an automated, flexible, and coherent digital engineering workflow. The simulation-to-Omniverse link is now organized around scenario-specific manifests that preserve the scene, terrain, route, actor, and DES inputs used for playback. CAD submission accepts several common source formats, distinguishes source engineering metadata from visualization metadata, and writes scaled physical information back to the SysML authoritative source of truth. The original ISRU DES is now transparent and configurable through the interface, while a role-based generic framework demonstrates how additional resource scenarios can be described and executed without creating a dedicated controller for every example.

The work also creates practical new capabilities for mission iteration. Users can save named scenario configurations, reload them, compare completed alternatives, and visualize each result through an independent Omniverse package. The extension provides terrain-aware placement, route-slope inspection, camera tools, and telemetry synchronized with the recorded DES. Together, these capabilities strengthen the digital thread between a stakeholder decision, the quantitative evaluation of that decision, and an immersive representation of its operational consequences.

The next priority is to mature the generic DES behavior language. Future work should add more complete resource inventories and queues, unit-aware validation, richer processing and storage behavior, stochastic events, detailed rover loading and charging operations, and a more complete power-network representation. These developments should be accompanied by verification cases and, where possible, validation against independently derived operational models. The scenario comparison interface should be extended to handle arbitrary resource names and more general measures of effectiveness rather than relying on ISRU-oriented result categories.

On the visualization side, the most important next step is to connect site selection in the interface to automatic extraction and preprocessing of the corresponding terrain patch from a larger lunar terrain database. This would make terrain consistency independent of manually prepared site assets. The Omniverse dashboard can then be expanded into a more resource-agnostic results interface and, eventually, connected to simulation progress while a run is executing rather than only to completed logs.

The CAD workflow provides a foundation for a curated component library. A future version can connect the existing repository CAD library to a broader catalog of lunar systems, certified metadata, material properties, and reusable subsystem models. Such a library would allow a user to select an existing component, inspect its provenance and capabilities, and incorporate it into a scenario without manually uploading a new asset each time.

Finally, true collaborative mission design requires an application architecture beyond the present static GitHub-based prototype. Named scenarios, saved configurations, and independent output packages provide a useful foundation for asynchronous collaboration. Real-time multi-user editing would require persistent backend storage, user accounts and access control, versioning and conflict resolution for shared scenarios, job management for concurrent simulations, and synchronized updates to all connected users. Addressing these requirements would transform ECLIPSE from a connected workflow into a collaborative mission-design application while preserving the explicit data contracts established in this work.
